% =====================================================================
% apendice.tex — Apêndice A: Roteiro de Demonstração
% =====================================================================

\chapter{Roteiro de Demonstração para a Banca}
\label{ap:roteiro}

Este apêndice consolida o roteiro de demonstração utilizado durante a
defesa do projeto, organizado em onze passos sequenciais que cobrem
desde a inicialização do sistema até a recalibração e o encerramento.
Cada passo é acompanhado da ação a ser executada e do resultado
esperado para fins de avaliação.

\begin{table}[!ht]
\centering
\caption{Roteiro de demonstração passo a passo.}
\label{tab:roteiro}
\small
\begin{tabularx}{\textwidth}{@{}c X X@{}}
\toprule
\textbf{\#} & \textbf{Ação} & \textbf{Resultado esperado} \\
\midrule
1  & Executar \texttt{python main.py}             & Abertura da GUI de configuração (CustomTkinter) \\
2  & Definir dimensões da mesa, modo de calibração (tampa/base) e mapa (GeoTIFF ou demo) & Janelas \texttt{Projecao\_Areia} e \texttt{Gabarito\_MDE} abertas \\
3  & Pressionar tecla \texttt{C}                  & Calibração executada (RANSAC + SVD + Gram-Schmidt + $T_{\text{shift}}$) \\
4  & Observar \texttt{Projecao\_Areia}            & Malha contínua: verde fora do cubo, azul no terço central (mapa Cubo Central) \\
5  & Observar \texttt{Gabarito\_MDE}              & Heatmap do MDE (Cubo Central, Morro Gaussiano ou GeoTIFF real) de referência \\
6  & Arrastar botão direito do mouse no terço central & Quadrados centrais transitam de azul $\to$ verde \\
7  & Arrastar botão esquerdo do mouse fora do cubo & Quadrados afetados transitam para azul, e voltam ao verde ao ajustar \\
8  & Continuar interação                          & Convergência da malha para totalmente verde \\
9  & Pressionar tecla \texttt{F}                  & Janela \texttt{Projecao\_Areia} entra em modo tela cheia \\
10 & Pressionar tecla \texttt{C} novamente        & Recalibração bem-sucedida (demonstra robustez) \\
11 & Pressionar tecla \texttt{Q} ou \texttt{ESC}  & Encerramento limpo, recursos liberados \\
\bottomrule
\end{tabularx}
\end{table}

\noindent
\textbf{Pontos-chave para destacar à banca:}
\begin{itemize}
    \item A interação com o \emph{mouse} demonstra em tempo real
          \emph{todo} o \textit{pipeline} matemático (discretização
          em malha $\to$ média espacial $\to$ comparação com MDE
          $\to$ projeção \textit{pinhole} de Tsai $\to$ rasterização
          por \texttt{cv2.fillPoly}), sem necessidade de
          hardware físico;
    \item A recalibração em tempo de execução comprova a robustez do
          mecanismo de \textit{fallback} e a estabilidade numérica do
          RANSAC seguido de SVD, mesmo após sucessivas modificações na
          areia virtual;
    \item A taxa de quadros (FPS) exibida no canto superior esquerdo
          da janela de projeção mantém-se consistentemente acima de
          30~fps, validando o requisito de operação em tempo real.
\end{itemize}

\section*{Comandos auxiliares}

Para executar a suíte de testes unitários:

\begin{lstlisting}
python -m unittest test_motor_caixao -v
\end{lstlisting}

\noindent Saída esperada: \texttt{88 passed}, com Open3D instalado.
Na ausência dessa biblioteca, o teste \texttt{TestLeituraRGBD} é
ignorado automaticamente, e a suíte reporta \texttt{87 passed} e
\texttt{1 skipped}.

Para forçar o modo simulação mesmo com Kinect conectado, editar no
topo do arquivo \texttt{main.py}:

\begin{lstlisting}
FORCAR_SIMULACAO = True
\end{lstlisting}
